\documentclass{article}

\usepackage[margin=1in]{geometry}
\usepackage{amsmath,amsthm,amssymb}
\usepackage[spanish,es-nodecimaldot]{babel}
\usepackage{enumitem}

\theoremstyle{definition}
\newtheorem{definition}{Definición}[section]

% Number sets
\newcommand{\R}{\mathbb{R}}
\newcommand{\Z}{\mathbb{Z}}
\newcommand{\N}{\mathbb{N}}
\newcommand{\Q}{\mathbb{Q}}
\newcommand{\C}{\mathbb{C}}
\newcommand{\K}{\mathbb{K}}

% Functional analysis operators
\newcommand{\norm}[1]{\left\lVert #1 \right\rVert}       % Norm
\newcommand{\seminorm}[1]{\left\lVert #1 \right\rVert}   % Seminorm (same symbol)
\newcommand{\cl}[1]{\overline{#1}}                       % Closure
\newcommand{\Int}[1]{{#1}^\circ}                         % Interior
\newcommand{\bola}[2]{B(#1;#2)}                          % Ball centered at #1 with radius #2
\newcommand{\EN}[1][E]{(#1,\norm{\cdot})}                % Normed space (E, ||.||); use \EN or \EN[F]
\newcommand{\ip}[2]{\left\langle #1, #2 \right\rangle}   % Inner product
\newcommand{\LH}{\mathcal{L}(H)}                         % Bounded operators on H
\newcommand{\Ker}{\operatorname{ker}}                    % Kernel
\newcommand{\Ran}{\operatorname{R}}                      % Range
\newcommand{\Lat}{\operatorname{Lat}}                    % Invariant subspaces lattice
\newcommand{\subesp}{\sqsubseteq}                        % Closed subspace

% Useful shortcuts
\newcommand{\sii}{\iff}                                  % Si y sólo si
\newcommand{\imp}{\implies}                              % Implies
\newcommand{\setdiff}{\smallsetminus}                    % Set difference
\newcommand{\diam}{\text{diam}}

% Exercise environment
\newenvironment{ejercicio}[1]{\begin{trivlist}
\item[\hskip \labelsep {\bfseries Ejercicio}\hskip \labelsep {\bfseries #1.}]}{\end{trivlist}}

\setlist[enumerate,1]{label=\alph*), leftmargin=*, itemsep=2pt}
\setlist[enumerate,2]{label=\roman*., leftmargin=*, itemsep=2pt}

\begin{document}

\title{Análisis Funcional: Trabajo Práctico N°6 - 2026}
\author{Operadores en Espacios de Hilbert}
\date{}

\maketitle

\vspace{0.8cm}

\section{Operadores en Espacios de Hilbert}

\begin{ejercicio}{1}
	Sean \( H \) un espacio de Hilbert complejo y \( T \in \LH \) un operador normal. Demostrar:
	\begin{enumerate}
		\item \( \norm{T}^2 = \norm{T^2} \).
		\item \( \Ker(T) = \Ker(T^*) \).
		\item Si \( T^2 = T \) entonces \( T \) es autoadjunto.
		\item Si \( T^2 = 0 \) entonces \( T = 0 \).
	\end{enumerate}
\end{ejercicio}

\vspace{0.5cm}

\begin{ejercicio}{2}
	Sean \( H \) un espacio de Hilbert complejo y \( S, T \in \LH \). Demostrar:
	\begin{enumerate}
		\item \( T^*T \) y \( TT^* \) son operadores positivos.
		\item Si \( T \) es positivo, entonces \( S^*TS \) es positivo.
	\end{enumerate}
\end{ejercicio}

\vspace{0.5cm}

\begin{ejercicio}{3}
	Sean \( H \) un espacio de Hilbert complejo y \( T \in \LH \). Demostrar:
	\begin{enumerate}
		\item Si \( \ip{Tx}{x} = 0 \) para todo \( x \in H \) entonces \( T = 0 \).
		\item Encontrar un contraejemplo que muestre la falsedad del enunciado anterior si \( H \) es un espacio de Hilbert real.
		\item \( T \) es normal si, y sólo si, \( \norm{Tx} = \norm{T^*x} \) para todo \( x \in H \).
		\item Probar que si \( H \) es un espacio de Hilbert real y \( T \in \LH \) autoadjunto entonces vale a).
	\end{enumerate}
\end{ejercicio}

\vspace{0.5cm}

\begin{ejercicio}{4}
	Dado un espacio de Hilbert \( H \), sean \( \{T_n\}, T \in \LH \) operadores normales. Probar que si \( T_n(x) \xrightarrow[n \to \infty]{} T(x) \) para todo \( x \in H \), entonces \( T_n^*(x) \xrightarrow[n \to \infty]{} T^*(x) \) para todo \( x \in H \).
\end{ejercicio}

\vspace{0.5cm}

\begin{ejercicio}{5}
	Sean \( H \) un espacio de Hilbert y \( P : H \to H \) un proyector. Demostrar que son equivalentes:
	\begin{enumerate}
		\item \( P^* = P \).
		\item \( P \) es normal.
		\item \( \norm{P} = 1 \).
		\item \( \Ker(P) = \Ran(P)^{\perp} \).
	\end{enumerate}
\end{ejercicio}

\vspace{0.5cm}

\begin{ejercicio}{6}
	Sean \( a, b \in \R \) con \( a < b \). Consideramos el espacio de Hilbert \( L^2(a, b) \) sobre \( \R \) y definimos el funcional \( T : L^2(a, b) \to \R \) dado por
	\[ T(f) = \int_a^b f(x) \, dx \]
	para \( f \in L^2(a, b) \).
	\begin{enumerate}
		\item Probar que \( T \) es acotado, lineal y calcular \( \norm{T} \).
		\item De acuerdo al Teorema de Riesz existe \( g \in L^2(a, b) \) tal que \( T(f) = \ip{f}{g} \) para toda \( f \in L^2(a, b) \). Hallar tal \( g \) explícitamente y verificar que \( \norm{g}_2 = \norm{T} \).
	\end{enumerate}
\end{ejercicio}

\vspace{0.5cm}

\begin{ejercicio}{7}
	Sea \( H \) un espacio de Hilbert y \( A \in \LH \) un operador autoadjunto. Probar que \( 0 \le A \le 1 \) si y sólo si \( A^2 \le A \).
\end{ejercicio}

\vspace{0.5cm}

\begin{ejercicio}{8}
	Sea \( T \in \mathcal{L}(L^2(0,1)) \) definido por \( (Tf)(x) = xf(x) \).
	\begin{enumerate}
		\item Probar que \( T \ge 0 \).
		\item Calcular \( T^{1/n} \) donde \( n \in \N \).
	\end{enumerate}
\end{ejercicio}

\vspace{0.5cm}

\begin{ejercicio}{9}
	Sea \( H \) un espacio de Hilbert. Sea \( A \in \LH^+ \) y tomemos su raíz cuadrada \( A^{1/2} \in \LH^+ \). Probar que
	\[ \Ker(A) = \Ker(A^{1/2}) \quad \text{pero} \quad \Ran(A) \subseteq \Ran(A^{1/2}) \subseteq \cl{\Ran(A)}. \]
	Mostrar además que si \( \Ran(A) \not\subesp H \) entonces las inclusiones son estrictas.
\end{ejercicio}

\vspace{0.5cm}

\begin{ejercicio}{10}
	Sea \( H \) un espacio de Hilbert y sea \( T \in \LH^+ \). Probar que \( T^2 \le \norm{T}\, T \).
\end{ejercicio}

\vspace{0.5cm}

\begin{ejercicio}{11}
	Probar que \( U \) es una isometría parcial si y sólo si \( U = UU^*U \).
\end{ejercicio}

\vspace{0.5cm}

\begin{ejercicio}{12}
	Hallar la Descomposición Polar de los siguientes operadores:
	\begin{enumerate}
		\item El operador de multiplicación \( M_\varphi : L^2([0,1]) \to L^2([0,1]) \) dado por
		      \[ M_\varphi(f) = \varphi f, \]
		      donde \( \varphi \in L^\infty([0,1]) \).
		\item Los operadores Shifts: \( S \in \mathcal{L}(\ell^2) \) y \( T \in \mathcal{L}(\ell^2) \) los shifts a izquierda y derecha, respectivamente.
		\item El operador \( A(x_1, x_2, \ldots) = (0, \alpha_1 x_1, \alpha_2 x_2, \ldots) \) donde \( x \in \ell^2 \) y \( \sup_{n \in \N} |\alpha_n| < \infty \).
	\end{enumerate}
\end{ejercicio}

\vspace{0.5cm}

\begin{ejercicio}{13}
	Dado un espacio de Hilbert \( H \), sea \( T \in \LH \). Probar que:
	\begin{enumerate}
		\item Si un subespacio \( S \) es \( T \)-invariante, entonces \( S \in \Lat(T) \).
		\item Si ahora \( S \subesp H \) y consideramos \( P_S \) el proyector ortogonal o autoadjunto, entonces
		      \begin{enumerate}
			      \item \( S \in \Lat(T) \) si y sólo si \( TP_S = P_S T P_S \).
			      \item \( S \in \Lat_r(T) \) si y sólo si \( TP_S = P_S T \), i.e. si \( T \) y \( P_S \) conmutan.
		      \end{enumerate}
	\end{enumerate}
\end{ejercicio}

\end{document}
